\section{Background and Related Work}
\label{sec:background}

This section discusses recent work from the compiler community that is closest to the study. Each cited paper is described by its contribution, and the comparison with the present study is stated at the end of each subsection. Tools and benchmarks are cited where the methods use them.

\subsection{Phase ordering and pass interactions}

Choosing passes and their order is a search problem. Ashouri et al.\ present the Protean Compiler, which orders optimizations at fine granularity under the guidance of a learned predictor \cite{ashouri2026protean}. Han and Kim expand the phase-ordering search space with dormant information, and they prune it more conservatively than methods that group transformations by expert knowledge \cite{han2024exponentially}. Xie et al.\ introduce optimization interaction, in which several phases, among them inlining and loop unrolling, act on the same code, and they fuzz a JVM compiler to maximize such interactions \cite{xie2024validating}. Seeker et al.\ discover heuristics that are embedded in a compiler rather than exposed as flags, and they tune those heuristics automatically \cite{seeker2024revealing}. Lenfers et al.\ model schedules as graph-structured objects, so that one auto-scheduling framework can serve several compiler infrastructures \cite{lenfers2026schedgehammer}. VenkataKeerthy et al.\ connect model development in Python to an optimizing compiler, so that learned models can make decisions inside it \cite{venkatakeerthy2024the}. Cummins et al.\ train language models on LLVM IR and assembly and apply them to optimization tasks \cite{cummins2025llm}. Gassmann et al.\ evaluate LLVM optimizations for zero-knowledge virtual machines, whose costs differ from those of CPUs, and they show that a few zkVM-aware changes to LLVM passes help \cite{gassmann2026evaluating}. dos Santos et al.\ show that optimization flags change the error rate of GPU programs, so that the effects of optimization reach beyond speed \cite{santos2024assessing}.

Much of the work above searches over pass sequences or measures the effect of a single optimization configuration. The present study instead measures how strongly the effect of one pass depends on the passes that run before it, in instruction count, and how much of that dependence the timing data can resolve.

\subsection{Vectorization, loop transformations, code size, and compile time}

Alladi et al.\ enable compiler-driven vectorization of Transformer models for edge devices, including RISC-V systems. Where the existing auto-vectorizer fails, they add a reduction-detection pass and a transpose step \cite{alladi2026enabling}. Murtovi et al.\ pair control-flow unmerging with loop unrolling on GPUs, so that the combination unlocks follow-on optimizations that neither technique reaches alone \cite{murtovi2024enhancing}. Zhi et al.\ retrieve demonstrations for a language model that rewrites loops in static control parts, and they check each candidate for equivalence with mutation, coverage, and differential testing \cite{zhi2026looprag}. H.\ Kwon et al.\ propose a chain of verification for language-model-generated transformations. Cheap static screens run first, symbolic equivalence checks based on Alive2 \cite{lopes2021alive2} run next, and run-time checks cover what cannot be proved statically. The chain widens vectorization coverage in LLVM \cite{kwon2026compiler}. Italiano and Cummins use a language model together with differential testing to find missed code-size optimizations in C and C++ compilers \cite{italiano2025finding}. Yuan et al.\ shrink linked binaries by outlining repeated instruction sequences, and they use runtime profiles to balance size against speed \cite{yuan2025post}. Saito et al.\ merge three or more similar functions at once, using multiple sequence alignment, where earlier approaches merged only pairs \cite{saito2025multiple}. Cugurovic et al.\ classify methods as cold or warm with a learned model, and they send cold methods to size-reducing transformations and warm methods to speed-oriented ones \cite{cugurovic2026graalmhc}. Mu et al.\ use sensor specifications to guide LLVM optimizations \cite{mu2024cosense}. Schwarz et al.\ describe TPDE, a back end that analyzes the code once and then generates it in a single sweep \cite{schwarz2026tpde}. Engelke and Schwarz compare how quickly several code generators compile query code \cite{engelke2024compile}. Titzer analyzes single-pass WebAssembly compilers and their trade-offs \cite{titzer2024whose}. Han et al.\ make the compiler stateful, so that incremental builds skip passes that were dormant in earlier builds \cite{han2024enabling}.

These studies report compile time, code size, or run time for their chosen configurations. The present study instead measures how a pair of passes interacts in instruction count, and it uses timing as a second, noisier measurement of the same interaction.

\subsection{Verification, miscompilations, and testing}

Fehr et al.\ write MLIR verification dialects that state what operations mean, so that optimizers and verifiers can rely on a shared semantics, and they report defects that they found in upstream MLIR \cite{fehr2025first}. Xu et al.\ couple a language model with a formal verifier in a feedback loop to find peephole optimizations that LLVM misses \cite{xu2026lpo}. Popescu and Lopes measure how much performance undefined-behavior-based optimizations deliver across processor families, and they find only slight end-to-end gains in LLVM \cite{popescu2025exploiting}. Mao et al.\ apply custom LLVM passes and bytecode rewriting to eBPF programs before the kernel verifier runs, and they confirm that the verifier accepts the tuned programs \cite{mao2024merlin}. Fan and Regehr integrate mutation-based fuzzing with Alive2 inside the LLVM middle end, so that mutation, optimization, and verification run in one process \cite{fan2024high}. J.\ Kwon et al.\ combine optimization-directed grey-box fuzzing with translation validation, to find incorrect optimizations in compilers that change continuously \cite{kwon2025optimization}. Lobo et al.\ set out a roadmap to remove undef from LLVM IR, because the interaction of undef with ordinary rewrites makes simple algebraic transformations unsound, and the semantics of undefined and poison values are subtle \cite{lee2017undefined,lobo2026towards}. Yu et al.\ present Ratte, which uses composable semantics for MLIR dialects to build modular fuzzers that find miscompilations \cite{yu2025ratte}.

Most of the work above targets individual optimizations or the compiler as a whole. The checker used here runs on every sweep, one module at a time, and it covers integer code only, as Section~\ref{sec:methods} states.
